\begin{afterword}
\summaryhead

The post uses the riddle ``If a tree falls in a forest, and no one hears it, does it make a sound?'' to stage what it calls the Standard Dispute. Albert says the tree makes a sound, because it makes vibrations in the air. Barry says it does not, because no one hears anything. They do not disagree about anything that happens in the forest, yet the argument continues. It moves to ``by definition,'' to a dictionary that lists both senses of ``sound,'' to accusations of breaking ``common usage,'' and to insults.

The author notes that the two agree on virtually every fact and still feel no agreement. In a comic scene, a time traveler proposes a remedy. The two can describe the event in unambiguous lower-level terms, such as acoustic vibrations and auditory experiences. Or each can adopt a new word, ``alberzle'' and ``bargulum,'' for what he meant by ``sound.'' Either way they should keep track of a testable proposition. The post ends by saying that this settles a substantial fraction of disputes over categories, though not every one.

\responsehead

The post's claim holds within its example. Albert and Barry agree on the vibrations and on the absence of a listener, and each side knows what the other means by ``sound.'' Once the two senses are separated, nothing is left to dispute. What the post contributes is the scene itself. It traces how a dispute of this kind moves from the forest to the word, then to the dictionary, then to ``common usage,'' and then to loss of face. It also observes that agreement on the facts brings no feeling of agreement.

The diagnosis and the remedy are old, and the post does not claim otherwise. The riddle was answered by a definition in its earliest known printing, in a magazine of 1883. William James told the same story in 1907, about a squirrel on a tree trunk (his general rule is quoted in the notes on ``Disguised Queries''). He separated two senses of ``going round,'' and he reported that one or two of the disputants insisted on ``plain honest English.'' The sentence ``Dictionary editors are historians of usage, not legislators of language'' closely echoes a sentence of S.~I. Hayakawa's from 1949, ``The writer of a dictionary is a historian, not a lawgiver''. Later philosophy made the renaming step formal: David Chalmers's ``subscript gambit'' of 2011 is the post's renaming of the two senses as ``alberzle'' and ``bargulum.'' These sources support the post; a reader should know that the ideas have a history.

How often the remedy works is asserted, not shown. ``A substantial fraction'' of disputes over categories are said to be settled this way, with no count and no example beyond the tree. The limit is real, and the next day's post names the disputes the remedy does not settle.

In short: a clear illustration of an old diagnosis, which William James gave with a squirrel in 1907, with a remedy that settles its own example, and a claim that it settles ``a substantial fraction'' of disputes that is made without evidence.
\end{afterword}
